%\documentclass[prl,aps,epsf,multicol]{revtex4-1}
%\documentclass[aps,prl,preprint,superscriptaddress]{revtex4-1}
%\documentclass[prl,aps,superscriptaddress,twocolumn]{revtex4-1}
\documentclass[prb,preprintnumbers,superscriptaddress,showpacs,preprint]{revtex4-1}
\usepackage{graphicx}
\usepackage{rotating}
\usepackage{dcolumn}% Align table columns on decimal point
\usepackage{bm}% bold math
\usepackage{amsmath,amssymb,amsfonts,textcomp}
\usepackage{multirow}
%\usepackage[hypertex]{hyperref}
\usepackage{color}
\usepackage{array}
\usepackage{hhline}
\usepackage[normalem]{ulem}
\usepackage{booktabs}
\usepackage{amsmath}    % need for subequations
\usepackage{verbatim}   % useful for program listings
\usepackage{subfigure}  % use for side-by-side figures
\tolerance=10000
%
\begin{document}

\title{The dispersion relation of excitations in superfluid $^4$He\\
SUPPLEMENTAL MATERIAL} 


\author{H. Godfrin}
 \email[Corresponding author: ]{henri.godfrin@neel.cnrs.fr}
  \affiliation{Univ. Grenoble Alpes, CNRS, Grenoble INP{\,$^\dagger$}, Institut N\'eel, 38000 Grenoble, France}
 \thanks{Institute of Engineering Univ. Grenoble Alpes}
\author{K. Beauvois} 
  \affiliation{Univ. Grenoble Alpes, CNRS, Grenoble INP{\,$^\dagger$}, Institut N\'eel, 38000 Grenoble, France}
	\thanks{Institute of Engineering Univ. Grenoble Alpes}
	  \affiliation{Institut Laue-Langevin, CS 20156, 38042 Grenoble Cedex 9, France}
\author{A. Sultan}	
  \affiliation{Univ. Grenoble Alpes, CNRS, Grenoble INP{\,$^\dagger$}, Institut N\'eel, 38000 Grenoble, France}
	\thanks{Institute of Engineering Univ. Grenoble Alpes}
	  \affiliation{Institut Laue-Langevin, CS 20156, 38042 Grenoble Cedex 9, France}
\author{E.~Krotscheck}
   \affiliation{Department of Physics, University at Buffalo, SUNY Buffalo NY 14260, USA}
   \affiliation{Institute for Theoretical Physics, Johannes Kepler University, A 4040 Linz, Austria}
\author{J. Dawidowski}
   \affiliation{Comisi\'on Nacional de Energ\'ia At\'omica and CONICET, Centro At\'omico Bariloche, (8400) San Carlos de Bariloche,
R\'io Negro, Argentina}
\author{B. F\aa k} 
   \affiliation{Institut Laue-Langevin, CS 20156, 38042 Grenoble Cedex 9, France}
\author{J. Ollivier}
   \affiliation{Institut Laue-Langevin, CS 20156, 38042 Grenoble Cedex 9, France}

\date{\today}


\maketitle 

\section{Introduction} 
We provide in this Supplemental Material additional information on:
\begin{itemize}
\item The analogy between the maxon wave-vector, and the Debye wave-vector of a periodic solid.
\item The extension of Landau's analytical model, developed here to describe quantitatively the neutron data. 
\item The appropriate phonon and roton parameters, needed to calculate analytically the low temperature properties of superfluid $^4$He. 
\item Data tables of the thermodynamical properties of superfluid $^4$He at saturated vapor pressure (SVP), calculated using the measured dispersion curve:
\begin{itemize}
\item TABLE I: Specific heat (total, phonon, and roton contributions), and total molar entropy, as a function of temperature.
\item TABLE II: Upper bound on the uncertainties of the specific heat.
\item TABLE III: Total density, normal fluid densities (total, phonon and roton
contributions), and same quantities divided by the total density, as a function of temperature.
\end{itemize}
\item Detailed data files of the thermodynamical properties of superfluid $^4$He at SVP. 
Data are shown for $T<$1.3\,K: above this temperature, the dispersion relation depends on temperature (see manuscript).
\begin{itemize}
\item Dispersion relation of superfluid $^4$He at SVP and very low temperatures.
\item The dispersion relation at several pressures, essentially covering the whole pressure range: P=0, 0.51, 1.02, 2.01, 5.01, 10.01, 24.08 bar. 
\item Detailed values of the specific heat and entropy vs. temperature at SVP.
\item Number density of roton and phonon excitations vs. temperature, at SVP.
\end{itemize}

\end{itemize}

\newpage
\section{Comparing maxon and Debye wave-vectors}
The following section illustrates the connection between the maxon wave-vector $k_M$ and the Debye wave-vector $k_D$ of a periodic system. 
The Debye wave-vector k$_D$=(6$\pi^2$n)$^{1/3}$ is defined by assigning one degree of freedom per atom to the longitudinal acoustic phonon mode. The density of states is integrated with an upper limit k$_{D}$ obtained from  
$N_a=(V/(2\pi^2))\int^{k_{D}}_{0}k^2dk$, 
where $N_a$ is Avogadro's number and V the molar volume (see manuscript). 

\begin{figure}[htbp]
	\centering
		\includegraphics[width=0.7\columnwidth]{kDebye-kMaxon.eps}
	\caption{The ratio of the maxon wave-vector and the Debye wave-vector calculated for the longitudinal phonon of a periodic system, as a function of density.}
	\label{fig:kDebye}
\end{figure}

The Debye model \cite{FetterWalecka} usually describes a solid, where the wave-vector is limited by the inverse of the lattice spacing. In a liquid, such a limitation does not exist. The hard core of the helium atoms, however, introduces a characteristic length and solid-like properties, like the very existence of roton excitations \cite{SmithSolid,NozSolid}. The Debye wave-vector is quantitatively similar to the maxon wave-vector (the ratio is very close to one), the maxon being analogous to a zone-boundary phonon. At high pressures, one can expect this analogy to work even better, and k$_M$ closely follows, indeed, the density dependence of k$_D$, as seen in Fig. \ref{fig:kDebye}.

\newpage
\section{Extension of Landau's analytical model}
Analytical expressions for the thermodynamical properties of superfluid $^4$He are given below.
They are particularly useful at temperatures below 0.7\,K, in particular when phonon and roton contributions 
must be considered separately, like in flow experiments, for instance. 

\subsection{Phonons}
The phonon dispersion relation is described here by the series expansion up to ${\cal O}\left(k^7\right)$ 
 
\begin{equation}
\epsilon^{Ph}(k)=\hbar\omega=\hbar c k(1+ \alpha_1 k +\alpha_2 k^2 + \alpha_3 k^3 + \alpha_4 k^4+ \alpha_5 k^5+ \alpha_6 k^6) 
\end{equation}

The corresponding inverse series describing the single-valued (phonon) branch of the spectrum is:

$k=a_1 \omega + a_2 \omega^2 + a_3 \omega^3 + a_4 \omega^4 + a_5 \omega^5 + a_6 \omega^6 + a_7 \omega^7 +{\cal O}\left(\omega^8\right) $

where the coefficients are 
\small
\begin{align*}
&a_1=\frac{1}{c} \nonumber\\
&a_2=-\frac{\alpha_1}{c^2} \nonumber\\
&a_3=\frac{2 \alpha_1^2-\alpha_2}{c^3}; \nonumber\\
&a_4=\frac{-5 \alpha_1^3+5 \alpha_1 \alpha_2-\alpha_3}{c^4} \nonumber\\
&a_5=\frac{14 \alpha_1^4-21 \alpha_1^2 \alpha_2+6 \alpha_1 \alpha_3+3 \alpha_2^2-\alpha_4}{c^5} \nonumber\\
&a_6=\frac{-42 \alpha_1^5+84 \alpha_1^3 \alpha_2-28 \alpha_1^2 \alpha_3-28 \alpha_1 \alpha_2^2+7 \alpha_1 \alpha_4+7 \alpha_2 \alpha_3-\alpha_5}{c^6} \nonumber\\
&a_7=\frac{132 \alpha_1^6-330 \alpha_1^4 \alpha_2+120 \alpha_1^3 \alpha_3+180 \alpha_1^2 \alpha_2^2-36 \alpha_1^2 \alpha_4-72 \alpha_1 \alpha_2 \alpha_3+8 \alpha_1 \alpha_5-12 \alpha_2^3+8 \alpha_2 \alpha_4+4 \alpha_3^2-\alpha_6}{c^7}\nonumber\\
\end{align*}

\normalsize
The density of states of an isotropic 3-dimensional phonon system is

\begin{equation}
D(\omega) d\omega=\left(\frac{1}{2 \pi^2} \right) k^2 V dk 
\label{dos}
\end{equation}

and hence

$D(\omega)=d_2 \omega^2 + d_3 \omega^3 + d_4 \omega^4 + d_5 \omega^5 + d_6 \omega^6 + d_7 \omega^7 + d_8 \omega^8  +{\cal O}\left(\omega^9\right) $, \\

with

\small
$d_2= \frac{V}{2 \pi ^2 c^3}$\\

$d_3= - \frac{2 \text{$\alpha_1$} V}{\pi ^2 c^4} $\\

$d_4= \frac{5 V \left(2 \text{$\alpha_1$}^2-\text{$\alpha_2$}\right)}{2 \pi ^2 c^5}+\frac{5 \text{$\alpha_1$}^2 V}{2 \pi ^2 c^5}$\\

$d_5= \frac{3 V \left(-5 \text{$\alpha_1$}^3+5 \text{$\alpha_1$} \text{$\alpha_2$}-\text{$\alpha_3$}\right)}{\pi ^2 c^6}-\frac{\text{$\alpha_1$}^3 V}{\pi ^2 c^6}-\frac{6 \text{$\alpha_1$} V \left(2 \text{$\alpha_1$}^2-\text{$\alpha_2$}\right)}{\pi ^2 c^6}$\\

$d_6= -\frac{7 \text{$\alpha_1$} V \left(-5 \text{$\alpha_1$}^3+5 \text{$\alpha_1$} \text{$\alpha_2$}-\text{$\alpha_3$}\right)}{\pi ^2 c^7}+\frac{7 \text{$\alpha_1$}^2 V \left(2 \text{$\alpha_1$}^2-\text{$\alpha_2$}\right)}{2 \pi ^2 c^7}+\frac{7 V \left(2 \text{$\alpha_1$}^2-\text{$\alpha_2$}\right)^2}{2 \pi ^2 c^7}+\frac{7 V \left(14 \text{$\alpha_1$}^4-21 \text{$\alpha_1$}^2 \text{$\alpha_2$}+6 \text{$\alpha_1$} \text{$\alpha_3$}+3 \text{$\alpha_2$}^2-\text{$\alpha_4$}\right)}{2 \pi ^2 c^7}$\\

$d_7=-\frac{4 \text{$\alpha_1$} V \left(2 \text{$\alpha_1$}^2-\text{$\alpha_2$}\right)^2}{\pi ^2 c^8}-\frac{8 \text{$\alpha_1$} V \left(14 \text{$\alpha_1$}^4-21 \text{$\alpha_1$}^2 \text{$\alpha_2$}+6 \text{$\alpha_1$} \text{$\alpha_3$}+3 \text{$\alpha_2$}^2-\text{$\alpha_4$}\right)}{\pi ^2 c^8} + \frac{4 \text{$\alpha_1$}^2 V \left(-5 \text{$\alpha_1$}^3+5 \text{$\alpha_1$} \text{$\alpha_2$}-\text{$\alpha_3$}\right)}{\pi ^2 c^8}$\\

$+\frac{8 V \left(2 \text{$\alpha_1$}^2-\text{$\alpha_2$}\right) \left(-5 \text{$\alpha_1$}^3+5 \text{$\alpha_1$} \text{$\alpha_2$}-\text{$\alpha_3$}\right)}{\pi ^2 c^8}
+\frac{4 V \left(-42 \text{$\alpha_1$}^5+84 \text{$\alpha_1$}^3 \text{$\alpha_2$}-28 \text{$\alpha_1$}^2 \text{$\alpha_3$}-28 \text{$\alpha_1$} \text{$\alpha_2$}^2+7 \text{$\alpha_1$} \text{$\alpha_4$}+7 \text{$\alpha_2$} \text{$\alpha_3$}-\text{$\alpha_5$}\right)}{\pi ^2 c^8} $\\

$d_8=\frac{9 V \left(-5 \text{$\alpha_1$}^3+5 \text{$\alpha_1$} \text{$\alpha_2$}-\text{$\alpha_3$}\right)^2}{2 \pi ^2 c^9} + \frac{3 V \left(2 \text{$\alpha_1$}^2-\text{$\alpha_2$}\right)^3}{2 \pi ^2 c^9} + \frac{9 V \left(2 \text{$\alpha_1$}^2-\text{$\alpha_2$}\right) \left(14 \text{$\alpha_1$}^4-21 \text{$\alpha_1$}^2 \text{$\alpha_2$}+6 \text{$\alpha_1$} \text{$\alpha_3$}+3 \text{$\alpha_2$}^2-\text{$\alpha_4$}\right)}{\pi ^2 c^9}$\\

$+ \frac{9 \text{$\alpha_1$}^2 V \left(14 \text{$\alpha_1$}^4-21 \text{$\alpha_1$}^2 \text{$\alpha_2$}+6 \text{$\alpha_1$} \text{$\alpha_3$}+3 \text{$\alpha_2$}^2-\text{$\alpha_4$}\right)}{2 \pi ^2 c^9} 
- \frac{9 \text{$\alpha_1$} V \left(2 \text{$\alpha_1$}^2-\text{$\alpha_2$}\right) \left(-5 \text{$\alpha_1$}^3+5 \text{$\alpha_1$} \text{$\alpha_2$}-\text{$\alpha_3$}\right)}{\pi ^2 c^9}$\\

$ - \frac{9 \text{$\alpha_1$} V \left(-42 \text{$\alpha_1$}^5+84 \text{$\alpha_1$}^3 \text{$\alpha_2$}-28 \text{$\alpha_1$}^2 \text{$\alpha_3$}-28 \text{$\alpha_1$} \text{$\alpha_2$}^2+7 \text{$\alpha_1$} \text{$\alpha_4$}+7 \text{$\alpha_2$} \text{$\alpha_3$}-\text{$\alpha_5$}\right)}{\pi ^2 c^9}$\\

$ + \frac{9 V \left(132 \text{$\alpha_1$}^6-330 \text{$\alpha_1$}^4 \text{$\alpha_2$}+120 \text{$\alpha_1$}^3 \text{$\alpha_3$}+180 \text{$\alpha_1$}^2 \text{$\alpha_2$}^2-36 \text{$\alpha_1$}^2 \text{$\alpha_4$}-72 \text{$\alpha_1$} \text{$\alpha_2$} \text{$\alpha_3$}+8 \text{$\alpha_1$} \text{$\alpha_5$}-12 \text{$\alpha_2$}^3+8 \text{$\alpha_2$} \text{$\alpha_4$}+4 \text{$\alpha_3$}^2-\text{$\alpha_6$}\right)}{2 \pi ^2 c^9}$\\
\normalsize

The thermodynamical system consists of a fixed number of atoms N in a constant volume V, kept at a temperature T. 
The appropriate thermodynamic potential is the Helmholz free energy F=E-TS, where E is the internal energy of the system, and S its entropy. 
Thermodynamical properties can be calculated as a function of temperature by elementary statistical physics using the phonon variables of the second quantization Hamiltonian \cite{FetterWalecka}. 
The number of phonons is not conserved; since it is not an independent variable, it is determined by minimization of the free energy F with respect to the phonon number, immediately showing, according to the definition of the chemical potential, that the phonon chemical potentials must vanish: $\mu_i^{Ph}$=0 (see the discussion in Ref. \onlinecite{FetterWalecka}).   
The internal energy E in the presence of density excitations (phonons, rotons...) is obtained using the dispersion relation $\epsilon(k)$,  
the Bose-Einstein distribution function $n(\epsilon(k))$, and the density of states (Eq. \ref{dos}): 

\begin{equation}
	\label{equ:energy1}
E = E_0 + (k_B V/2\pi^2)\int_{0}^{\infty}\frac{\epsilon(k)}{e^{\frac{\epsilon(k)}{k_B T}}-1}k^2dk
\end{equation}
where $E_0$ is the (constant) phonon zero-point energy.

The upper limit of the integral is taken as infinity, the exponential factor making  this choice possible in the very low temperature limit. 
The integral can be performed over energies instead of wave-vectors, using Eq. \ref{dos}: 
\begin{equation}
	\label{equ:energy2}
 E = E_0 + \int_{0}^{\infty} D(\omega) \hbar \omega  \left(e^{ \frac{\hbar \omega}{k_B T}}-1\right)^{-1} d\omega
\end{equation}

which yields, using the series expansion for the density of states calculated above:  

\begin{equation}
E^{Ph}=E_0 + \frac{\textbf{A}}{4} T^4 + \frac{\textbf{B}}{5} T^5 + \frac{\textbf{C}}{6} T^6 + \frac{\textbf{D}}{7} T^7+\frac{\textbf{E}}{8} T^8 + \frac{\textbf{K}}{9} T^9+\frac{\textbf{L}}{10} T^{10}
\end{equation}

where 

$\textbf{A}=\frac{2 \pi ^2 k_B^4 V}{15 c^3 \hbar ^3} $\\

$\textbf{B}=-\frac{240 \left(\alpha_1 k_B^5 V \zeta (5)\right)}{\pi ^2 c^4 \hbar ^4} $\\

$\textbf{C}=\frac{40 \pi ^4 k_B^6 V \left(3 \alpha_1^2-\alpha_2\right)}{21 c^5 \hbar ^5} $\\

$\textbf{D}=-\frac{5040 \left(k_B^7 V \zeta (7) \left(28 \alpha_1^3-21 \alpha_1 \alpha_2+3 \alpha_3\right)\right)}{\pi ^2 c^6 \hbar ^6} $\\

$\textbf{E}=\frac{224 \pi ^6 k_B^8 V \left(30 \alpha_1^4-36 \alpha_1^2 \alpha_2+8 \alpha_1 \alpha_3+4 \alpha_2^2-\alpha_4\right)}{15 c^7 \hbar ^7}$\\

$\textbf{K}=\frac{1451520 \left(k_B^9 V \zeta (9) \left(99 \alpha_1^5-165 \alpha_1^3 \alpha_2+45 \alpha_1^2 \alpha_3+45 \alpha_1 \alpha_2^2-9 \alpha_1 \alpha_4-9 \alpha_2 \alpha_3+\alpha_5\right)\right)}{\pi ^2 c^8 \hbar ^8} $\\

$\textbf{L}=\frac{640 \pi ^8 k_B^{10} V \left(1001 \alpha_1^6-2145 \alpha_1^4 \alpha_2+660 \alpha_1^3 \alpha_3+165 \alpha_1^2 \left(6 \alpha_2^2-\alpha_4\right)+30 \alpha_1 (\alpha_5-11 \alpha_2 \alpha_3)-55 \alpha_2^3+30 \alpha_2 \alpha_4+15 \alpha_3^2-3 \alpha_6\right)}{11 c^9 \hbar ^9} $\\

where $\zeta(s)$ is the Riemann zeta function.\\

The specific heat at constant volume is then calculated with the expression $C_V = \left(\frac{\partial E}{\partial T}\right)_{V}$

\begin{equation}
C_{V}^{Ph}=\textbf{A} T^3 + \textbf{B} T^4+  \textbf{C} T^5+ \textbf{D} T^6+ \textbf{E} T^7+ \textbf{K} T^8+ \textbf{L} T^9
\end{equation}

The entropy can be obtained by integrating $dS=\frac{C_V}{T} dT$, i.e., 
$S(T)=\int_0^T \frac{C_V(T')}{T'} dT'$

\begin{equation}
S^{Ph}=\frac{\textbf{A}}{3} T^3+\frac{\textbf{B}}{4} T^4+\frac{\textbf{C}}{5} T^5+\frac{\textbf{D}}{6} T^6+\frac{\textbf{E}}{7} T^7+\frac{\textbf{K}}{8} T^8+\frac{\textbf{L}}{9} T^9
\end{equation}

and the Helmholz free energy is $F=E-T S$,

\begin{equation}
F^{Ph}=E_0 - \frac{\textbf{A}}{12} T^4 - \frac{\textbf{B}}{20} T^5 - \frac{\textbf{C}}{30} T^6 - \frac{\textbf{D}}{42} T^7-\frac{\textbf{E}}{56} T^8 - \frac{\textbf{K}}{72} T^9-\frac{\textbf{L}}{90} T^{10}
\end{equation}
\\
It has been assumed in the manuscript that the $\alpha_1$ term of the dispersion relation is exactly zero. This assumption is supported by measurements \cite{Rugar1984} that placed an upper limit to this coefficient of 10$^{-3}$\,\AA. In addition, there are no theoretical indications that this coefficient could have a finite value. Note, however, that even a relatively small value of $\alpha_1$ would significantly affect the coefficients \textbf{B}, \textbf{C}, \textbf{D}, etc..., where $\alpha_1$ systematically appears multiplied by a large numerical factor.

The values of the phonon parameters at saturated vapor pressure obtained in this work are, when temperatures are expressed in kelvin and C$_V$ in J/(mol.K):

\noindent \textbf{A}=0.0831 \\
\textbf{B}=0\\
\textbf{C}=-0.0548\\
\textbf{D}=0.0653 \\
\textbf{E}=0.0603 \\
\textbf{K}=-0.262 \\ 
\textbf{L}=0.141\\

\subsection{Rotons}
The analytical calculation of the roton contribution to the thermodynamical properties has been formulated by Landau \cite{Landauroton,Landauroton2} in his 1947 paper. 
The excitation energies around the roton minimum are described by the expansion

\begin{equation}
	\epsilon_R(k)=\Delta_R+\frac{\hbar^2}{2m_4\mu_R}(k-k_R)^2+B_R(k-k_R)^3+C_R(k-k_R)^4
	\label{Eq:roton}
\end{equation}

where $\Delta_R$ is the roton energy, $k_R$ the roton wave-vector, and $\mu_R$ the roton effective mass. We have introduced here B$_R$ and C$_R$, two parameters that describe the deformation of the roton minimum with respect to the simple parabolic shape assumed by Landau. 

The statistical physics calculation \cite{FetterWalecka} of the thermodynamic functions follows the lines already discussed above for the small wave-vector phonons. The internal energy is calculated using Eq. \ref{equ:energy1}, integrating now around the roton wave-vector.
Landau obtains the expression \cite{Landauroton2} : 

\begin{equation}
C_V^R=C_{\textit{L}} e^{-\frac{\Delta_R}{k_B T}} \left(1 +\frac{k_B T}{\Delta_R} +\frac{3}{4} \left(\frac{k_B T}{\Delta_R}\right)^2\right)
\label{eq:CvRotonLandau}
\end{equation}

with  $C_{\textit{L}}=\frac{\Delta_R^2 k_R^2 V \sqrt{\mu_R m_4}}{\sqrt{2} \pi^{3/2} \sqrt{k_B} T^{3/2} \hbar}$\\

The analysis of the measured specific heats with Landau's analytical expression is therefore sensitive to a combination of several roton parameters (energy, wave-vector and effective mass). 

Landau also assumes that $\left|k-k_R\right|<<k_R$. 
The analytical calculation for a parabolic dispersion relation can be done, however, without this last assumption. 
We calculate here the average energy of the rotons E$_R$ and their specific heat C$_V^R$ with Eq. \ref{equ:energy1} using Eq. \ref{Eq:roton} with B$_R$=C$_R$=0, and including in the integrand three terms originating from the expansion of $(k-k_R)^2$. The calculation is performed by integration over wave-vectors. The first term gives Landau's expression,the second term vanishes by symmetry, and  an additional term is found:

\begin{equation}
C_V^{addit}=C_{\textit{L}} e^{-\frac{\Delta_R}{k_B T}}\left(\frac{T }{2\tau_R}\right)\left(1+\frac{3 k_B T}{\Delta_R}+\frac{15}{4}\left(\frac{k_B T}{\Delta_R}\right)^2\right)
\label{eq:CvRotonHG}
\end{equation}

where we introduced the parameter $\tau_R$ $=\frac{k_R^2 \hbar^2}{2 k_B \mu_R m_4}$,  
 which has the dimensions of temperature.
The additional term is rather similar to Landau's expression, but its magnitude is affected by an additional multiplicative coefficient T/(2$\tau_R$). Since $\tau_R$ $\approx$ 158\,K, this additional term is of the same magnitude as the third term in Eq. \ref{eq:CvRotonLandau} for temperatures above 1\,K, and they are both very small.
In fact, the integral appearing in the additional term has the form found by Landau in his original 1941 paper \cite{Landauroton}, where the roton was supposed to be at zero wave-vector. 

As shown in the manuscript, it is necessary to take into account the large asymmetry (cubic term) and the deformation of the minimum
(quartic term). Unfortunately, 
we have not found a simple analytical expression for the case where the coefficients B or C in Eq. \ref{eq:CvRotonHG} are finite. 
At the analytical level, we are therefore essentially left with Landau's formula, of fundamental interest, but of modest accuracy, and limited to temperatures well below 1\,K (see manuscript).  

The roton parameters at saturated vapor pressure are:

 \noindent $\Delta_R$= \textit{0.7418(10)}\,meV = 8.608\,K (from Ref. \onlinecite{stirling-83,GlydeBook,StirlingExeter})\\
 $k_R$=1.918(2)\, \AA$^{-1}$\\
 $\mu_R$=0.141(2)\\
 B$_R$=2.64(5)\, \AA$^{3}$, and  \\
 C$_R$=-10.5(5)\, \AA$^{4}$  \\

Graphs providing phonon and roton parameters at other pressures are given in the manuscript.




\newpage
\section{TABLES}
\label{tablesSuppMat}
\subsection{Specific heat and entropy at SVP: numerical calculation}
\begingroup
\squeezetable
% Please add the following required packages to your document preamble:
% \usepackage{booktabs}
\begin{table}[!h]
\begin{tabular}{@{}|c|l|l|l|l|l|@{}}
\hline \hline
\multicolumn{1}{|c|}{Temperature} & \multicolumn{1}{c|}{Cv  (total) } & \multicolumn{1}{c|}{Cv  (phonons) } & \multicolumn{1}{c|}{Cv  (rotons)} &  & \multicolumn{1}{c|}{S (total)} \\ \midrule \hline
\multicolumn{1}{|c|}{\textit{K}}  & \multicolumn{1}{c|}{J/(K.mol)}  & \multicolumn{1}{c|}{J/(K.mol)}    & \multicolumn{1}{c|}{J/(K.mol)}   &  & \multicolumn{1}{c|}{J/(K.mol)} \\ \midrule \hline 
0.01	&	8.309E-08	&	8.309E-08	&	0.000	&	&	2.763E-08	\\	\midrule
0.02	&	6.645E-07	&	6.645E-07	&	6.550E-182	&	&	2.214E-07	\\	\midrule	
0.03	&	2.242E-06	&	2.242E-06	&	7.090E-120	&	&	7.471E-07	\\	\midrule	
0.04	&	5.312E-06	&	5.312E-06	&	6.500E-89	&	&	1.771E-06	\\	\midrule
0.05	&	1.037E-05	&	1.037E-05	&	2.270E-70	&	&	3.459E-06	\\	\midrule
0.06	&	1.791E-05	&	1.791E-05	&	4.980E-58	&	&	5.974E-06	\\	\midrule
0.07	&	2.841E-05	&	2.841E-05	&	3.130E-49	&	&	9.482E-06	\\	\midrule
0.08	&	4.238E-05	&	4.238E-05	&	1.210E-42	&	&	1.415E-05	\\	\midrule
0.09	&	6.028E-05	&	6.028E-05	&	1.580E-37	&	&	2.013E-05	\\	\midrule
0.10	&	8.260E-05	&	8.260E-05	&	1.920E-33	&	&	2.760E-05	\\	\midrule
0.11	&	1.098E-04	&	1.098E-04	&	4.170E-30	&	&	3.670E-05	\\	\midrule
0.12	&	1.424E-04	&	1.424E-04	&	2.480E-27	&	&	4.762E-05	\\	\midrule
0.13	&	1.808E-04	&	1.808E-04	&	5.490E-25	&	&	6.049E-05	\\	\midrule
0.14	&	2.255E-04	&	2.255E-04	&	5.570E-23	&	&	7.550E-05	\\	\midrule
0.15	&	2.770E-04	&	2.770E-04	&	3.030E-21	&	&	9.276E-05	\\	\midrule
0.16	&	3.357E-04	&	3.357E-04	&	9.950E-20	&	&	1.124E-04	\\	\midrule
0.17	&	4.020E-04	&	4.020E-04	&	2.150E-18	&	&	1.347E-04	\\	\midrule
0.18	&	4.764E-04	&	4.764E-04	&	3.290E-17	&	&	1.598E-04	\\	\midrule
0.19	&	5.594E-04	&	5.594E-04	&	3.770E-16	&	&	1.877E-04	\\	\midrule
0.20	&	6.513E-04	&	6.513E-04	&	3.360E-15	&	&	2.187E-04	\\	\midrule
0.21	&	7.527E-04	&	7.527E-04	&	2.430E-14	&	&	2.529E-04	\\	\midrule
0.22	&	8.639E-04	&	8.639E-04	&	1.460E-13	&	&	2.904E-04	\\	\midrule
0.23	&	9.853E-04	&	9.853E-04	&	7.520E-13	&	&	3.315E-04	\\	\midrule
0.24	&	0.001117	&	0.001117	&	3.360E-12	&	&	3.762E-04	\\	\midrule
0.25	&	0.001261	&	0.001261	&	1.330E-11	&	&	4.247E-04	\\	\midrule
0.26	&	0.001415	&	0.001415	&	4.720E-11	&	&	4.771E-04	\\	\midrule
0.27	&	0.001582	&	0.001582	&	1.520E-10	&	&	5.336E-04	\\	\midrule
0.28	&	0.001761	&	0.001761	&	4.510E-10	&	&	5.943E-04	\\	\midrule
0.29	&	0.001953	&	0.001953	&	1.230E-09	&	&	6.594E-04	\\	\midrule
0.30	&	0.002157	&	0.002157	&	3.170E-09	&	&	7.290E-04	\\	\midrule
0.31	&	0.002376	&	0.002376	&	7.630E-09	&	&	8.033E-04	\\	\midrule
0.32	&	0.002608	&	0.002608	&	1.735E-08	&	&	8.823E-04	\\	\midrule
0.33	&	0.002854	&	0.002854	&	3.752E-08	&	&	9.663E-04	\\	\midrule
0.34	&	0.003116	&	0.003116	&	7.742E-08	&	&	0.001055	\\	\midrule
0.35	&	0.003392	&	0.003392	&	1.531E-07	&	&	0.001149	\\	\midrule
0.36	&	0.003684	&	0.003684	&	2.911E-07	&	&	0.001249	\\	\midrule
0.37	&	0.003992	&	0.003991	&	5.343E-07	&	&	0.001354	\\	\midrule
0.38	&	0.004316	&	0.004315	&	9.488E-07	&	&	0.001465	\\	\midrule
0.39	&	0.004657	&	0.004655	&	1.634E-06	&	&	0.001582	\\	\bottomrule	\hline 
\end{tabular}
\end{table}
\endgroup


\newpage
\begingroup
\squeezetable
% Please add the following required packages to your document preamble:
% \usepackage{booktabs}
\begin{table}[!h]
\begin{tabular}{@{}|c|l|l|l|l|l|@{}}
\hline 
\multicolumn{1}{|c|}{Temperature} & \multicolumn{1}{c|}{Cv  (total) } & \multicolumn{1}{c|}{Cv  (phonons) } & \multicolumn{1}{c|}{Cv  (rotons)} &  & \multicolumn{1}{c|}{S (total)} \\ \midrule \hline
\multicolumn{1}{|c|}{\textit{K}}  & \multicolumn{1}{c|}{J/(K.mol)}  & \multicolumn{1}{c|}{J/(K.mol)}    & \multicolumn{1}{c|}{J/(K.mol)}   &  & \multicolumn{1}{c|}{J/(K.mol)} \\ \midrule \hline 
0.40	&	0.005016	&	0.005013	&	2.738E-06	&	&	0.001704	\\	\midrule	
0.41	&	0.005392	&	0.005388	&	4.468E-06	&	&	0.001832	\\	\midrule	
0.42	&	0.005787	&	0.005780	&	7.119E-06	&	&	0.001967	\\	\midrule	
0.43	&	0.006202	&	0.006191	&	1.109E-05	&	&	0.002108	\\	\midrule	
0.44	&	0.006637	&	0.006620	&	1.692E-05	&	&	0.002255	\\	\midrule	
0.45	&	0.007094	&	0.007069	&	2.532E-05	&	&	0.002409	\\	\midrule	
0.46	&	0.007573	&	0.007536	&	3.721E-05	&	&	0.002570	\\	\midrule	
0.47	&	0.008077	&	0.008023	&	5.377E-05	&	&	0.002739	\\	\midrule	
0.48	&	0.008607	&	0.008531	&	7.647E-05	&	&	0.002914	\\	\midrule	
0.49	&	0.009166	&	0.009059	&	1.071E-04	&	&	0.003097	\\	\midrule	
0.50	&	0.009756	&	0.009607	&	1.480E-04	&	&	0.003288	\\	\midrule	
0.51	&	0.01038	&	0.01018	&	2.018E-04	&	&	0.003488	\\	\midrule	
0.52	&	0.01104	&	0.01077	&	2.717E-04	&	&	0.003696	\\	\midrule	
0.53	&	0.01175	&	0.01138	&	3.617E-04	&	&	0.003913	\\	\midrule	
0.54	&	0.01250	&	0.01202	&	4.761E-04	&	&	0.004139	\\	\midrule	
0.55	&	0.01330	&	0.01268	&	6.203E-04	&	&	0.004376	\\	\midrule	
0.56	&	0.01416	&	0.01336	&	8.003E-04	&	&	0.004623	\\	\midrule	
0.57	&	0.01509	&	0.01407	&	0.001022	&	&	0.004882	\\	\midrule	
0.58	&	0.01610	&	0.01480	&	0.001295	&	&	0.005153	\\	\midrule	
0.59	&	0.01719	&	0.01556	&	0.001627	&	&	0.005437	\\	\midrule	
0.60	&	0.01837	&	0.01634	&	0.002028	&	&	0.005736	\\	\midrule	
0.61	&	0.01966	&	0.01715	&	0.002509	&	&	0.006050	\\	\midrule	
0.62	&	0.02106	&	0.01798	&	0.003082	&	&	0.006381	\\	\midrule	
0.63	&	0.02260	&	0.01884	&	0.003759	&	&	0.006730	\\	\midrule	
0.64	&	0.02428	&	0.01973	&	0.004556	&	&	0.007099	\\	\midrule	
0.65	&	0.02613	&	0.02064	&	0.005487	&	&	0.007489	\\	\midrule	
0.66	&	0.02816	&	0.02159	&	0.006570	&	&	0.007903	\\	\midrule	
0.67	&	0.03038	&	0.02256	&	0.007823	&	&	0.008343	\\	\midrule	
0.68	&	0.03283	&	0.02356	&	0.009264	&	&	0.008811	\\	\midrule	
0.69	&	0.03551	&	0.02460	&	0.01091	&	&	0.009310	\\	\midrule	
0.70	&	0.03846	&	0.02566	&	0.01279	&	&	0.009841	\\	\midrule	
0.71	&	0.04169	&	0.02676	&	0.01493	&	&	0.01041	\\	\midrule	
0.72	&	0.04523	&	0.02788	&	0.01734	&	&	0.01101	\\	\midrule	
0.73	&	0.04911	&	0.02905	&	0.02006	&	&	0.01166	\\	\midrule	
0.74	&	0.05335	&	0.03024	&	0.02311	&	&	0.01236	\\	\midrule	
0.75	&	0.05784	&	0.03147	&	0.02651	&	&	0.01310	\\	\midrule	
0.76	&	0.06304	&	0.03273	&	0.03030	&	&	0.01390	\\	\midrule	
0.77	&	0.06855	&	0.03404	&	0.03451	&	&	0.01476	\\	\midrule	
0.78	&	0.07455	&	0.03537	&	0.03917	&	&	0.01569	\\	\midrule	
0.79	&	0.08106	&	0.03675	&	0.04430	&	&	0.01668	\\	\bottomrule	\hline
\end{tabular}
\end{table}
\endgroup


\newpage
\begingroup
\squeezetable
% Please add the following required packages to your document preamble:
% \usepackage{booktabs}
\begin{table}[!h]
\begin{tabular}{@{}|c|l|l|l|l|l|@{}}
\hline
\multicolumn{1}{|c|}{Temperature} & \multicolumn{1}{c|}{Cv  (total) } & \multicolumn{1}{c|}{Cv  (phonons) } & \multicolumn{1}{c|}{Cv  (rotons)} &  & \multicolumn{1}{c|}{S (total)} \\ \midrule \hline
\multicolumn{1}{|c|}{\textit{K}}  & \multicolumn{1}{c|}{J/(K.mol)}  & \multicolumn{1}{c|}{J/(K.mol)}    & \multicolumn{1}{c|}{J/(K.mol)}   &  & \multicolumn{1}{c|}{J/(K.mol)} \\ \midrule \hline 
0.80	&	0.08812	&	0.03816	&	0.04994	&	&	0.01774	\\	\midrule	
0.81	&	0.09576	&	0.03962	&	0.05613	&	&	0.01888	\\	\midrule	
0.82	&	0.1040	&	0.04112	&	0.06290	&	&	0.02011	\\	\midrule	
0.83	&	0.1129	&	0.04265	&	0.07028	&	&	0.02142	\\	\midrule	
0.84	&	0.1225	&	0.04423	&	0.07831	&	&	0.02283	\\	\midrule	
0.85	&	0.1329	&	0.04586	&	0.08702	&	&	0.02434	\\	\midrule	
0.86	&	0.1440	&	0.04753	&	0.09645	&	&	0.02596	\\	\midrule	
0.87	&	0.1559	&	0.04925	&	0.1066	&	&	0.02769	\\	\midrule	
0.88	&	0.1687	&	0.05101	&	0.1176	&	&	0.02955	\\	\midrule	
0.89	&	0.1823	&	0.05283	&	0.1294	&	&	0.03153	\\	\midrule	
0.90	&	0.1968	&	0.05469	&	0.1421	&	&	0.03364	\\	\midrule	
0.91	&	0.2124	&	0.05661	&	0.1557	&	&	0.03590	\\	\midrule	
0.92	&	0.2289	&	0.05857	&	0.1702	&	&	0.03831	\\	\midrule	
0.93	&	0.2464	&	0.06060	&	0.1858	&	&	0.04088	\\	\midrule	
0.94	&	0.2651	&	0.06268	&	0.2023	&	&	0.04362	\\	\midrule	
0.95	&	0.2848	&	0.06481	&	0.2200	&	&	0.04653	\\	\midrule	
0.96	&	0.3057	&	0.06701	&	0.2387	&	&	0.04962	\\	\midrule	
0.97	&	0.3279	&	0.06927	&	0.2585	&	&	0.05290	\\	\midrule	
0.98	&	0.3512	&	0.07158	&	0.2796	&	&	0.05638	\\	\midrule	
0.99	&	0.3759	&	0.07397	&	0.3018	&	&	0.06007	\\	\midrule	
1.00	&	0.4018	&	0.07641	&	0.3253	&	&	0.06397	\\	\midrule	
1.01	&	0.4291	&	0.07893	&	0.3501	&	&	0.06811	\\	\midrule	
1.02	&	0.4578	&	0.08151	&	0.3763	&	&	0.07248	\\	\midrule	
1.03	&	0.4880	&	0.08416	&	0.4038	&	&	0.07709	\\	\midrule	
1.04	&	0.5196	&	0.08689	&	0.4327	&	&	0.08195	\\	\midrule	
1.05	&	0.5527	&	0.08969	&	0.4630	&	&	0.08708	\\	\midrule	
1.06	&	0.5874	&	0.09256	&	0.4948	&	&	0.09249	\\	\midrule	
1.07	&	0.6237	&	0.09552	&	0.5281	&	&	0.09817	\\	\midrule	
1.08	&	0.6616	&	0.09855	&	0.5630	&	&	0.1041	\\	\midrule	
1.09	&	0.7012	&	0.1017	&	0.5995	&	&	0.1104	\\	\midrule	
1.10	&	0.7424	&	0.1049	&	0.6375	&	&	0.1170	\\	\midrule	
1.11	&	0.7854	&	0.1081	&	0.6773	&	&	0.1239	\\	\midrule	
1.12	&	0.8302	&	0.1115	&	0.7187	&	&	0.1311	\\	\midrule	
1.13	&	0.8768	&	0.1150	&	0.7618	&	&	0.1387	\\	\midrule	
1.14	&	0.9252	&	0.1185	&	0.8067	&	&	0.1466	\\	\midrule	
1.15	&	0.9755	&	0.1221	&	0.8533	&	&	0.1549	\\	\midrule	
1.16	&	1.028	&	0.1259	&	0.9018	&	&	0.1636	\\	\midrule	
1.17	&	1.082	&	0.1297	&	0.9521	&	&	0.1727	\\	\midrule	
1.18	&	1.138	&	0.1336	&	1.004	&	&	0.1821	\\	\midrule	
1.19	&	1.196	&	0.1376	&	1.058	&	&	0.1920	\\	\bottomrule	 \hline
\end{tabular}
\end{table}
\endgroup


\newpage
\begingroup
\squeezetable
% Please add the following required packages to your document preamble:
% \usepackage{booktabs}
\begin{table}[!h]
\begin{tabular}{@{}|c|l|l|l|l|l|@{}}
\hline 
\multicolumn{1}{|c|}{Temperature} & \multicolumn{1}{c|}{Cv  (total) } & \multicolumn{1}{c|}{Cv  (phonons) } & \multicolumn{1}{c|}{Cv  (rotons)} &  & \multicolumn{1}{c|}{S (total)} \\ \midrule \hline
\multicolumn{1}{|c|}{\textit{K}}  & \multicolumn{1}{c|}{J/(K.mol)}  & \multicolumn{1}{c|}{J/(K.mol)}    & \multicolumn{1}{c|}{J/(K.mol)}   &  & \multicolumn{1}{c|}{J/(K.mol)} \\ \midrule \hline 
1.20	&	1.256	&	0.1418	&	1.115	&	&	0.2022	\\	\midrule	
1.21	&	1.319	&	0.1460	&	1.173	&	&	0.2129	\\	\midrule	
1.22	&	1.383	&	0.1503	&	1.233	&	&	0.2240	\\	\midrule	
1.23	&	1.447	&	0.1547	&	1.295	&	&	0.2356	\\	\midrule	
1.24	&	1.518	&	0.1593	&	1.359	&	&	0.2476	\\	\midrule	
1.25	&	1.589	&	0.1639	&	1.425	&	&	0.2600	\\	\midrule	
1.26	&	1.662	&	0.1687	&	1.493	&	&	0.2730	\\	\midrule	
1.27	&	1.737	&	0.1736	&	1.564	&	&	0.2864	\\	\midrule	
1.28	&	1.815	&	0.1785	&	1.636	&	&	0.3004	\\	\midrule	
1.29	&	1.896	&	0.1836	&	1.711	&	&	0.3148	\\	\midrule	
1.30	&	1.977	&	0.1889	&	1.788	&	&	0.3297	\\	\bottomrule	
\hline \hline 
\end{tabular}
\caption{Specific heat (total, phonon, and roton contributions), and total molar entropy, at saturated vapor pressure, as a function of temperature, calculated numerically using the measured dispersion relation (this work).}
\label{tab:Cv+S}
\end{table}
\endgroup


\begingroup
\squeezetable
% Please add the following required packages to your document preamble:
% \usepackage{booktabs}
\begin{table}[!h]
\begin{tabular}{@{}|c|l|l|c|@{}}
\hline 
\multicolumn{1}{|c|}{Temperature} & \multicolumn{1}{c|}{Cv  (total) } & \multicolumn{1}{c|}{error Cv  } & \multicolumn{1}{c|}{rel. error Cv}  \\ \midrule \hline
\multicolumn{1}{|c|}{\textit{K}}  & \multicolumn{1}{c|}{J/(K.mol)}  & \multicolumn{1}{c|}{J/(K.mol)}    & \multicolumn{1}{c|}{\%}   \\ \midrule \hline 
0.05 &1.037E-5	&1.3E-8	&0.13	\\	\midrule
0.1	 &8.26E-5	  &1.1E-7	&0.13	\\	\midrule
0.2	 &6.513E-4	&1.1E-6	&0.17	\\	\midrule
0.3	 &0.002157	&5.7E-6	&0.26	\\	\midrule
0.4	 &0.005016	&1.8E-5	&0.36	\\	\midrule
0.5	 &0.009756	&3.7E-5	&0.38	\\	\midrule
0.6  &0.01837	  &9.2E-5	&0.50	\\	\midrule
0.7  &0.03846	  &3.3E-4	&0.86	\\	\midrule
0.8  &0.08812	  &9.9E-4	&1.1	\\	\midrule
0.9	 &0.1968	  &0.0023	&1.2	\\	\midrule
1.0  &0.4018	  &0.0045	&1.1	\\	\midrule
1.1	 &0.7424	  &0.0083	&1.1	\\	\midrule
1.2	 &1.256	    &0.013	&1.0	\\	\midrule 
1.3	 &1.977	    &0.021	&1.1	\\	\bottomrule
\hline \hline 
\end{tabular}
\caption{Upper bound on the uncertainties of the specific heat at  saturated vapor pressure, calculated numerically as a function of temperature using the measured dispersion relation. The uncertainty has been calculated assuming a systematic shift of the dispersion relation, in the same direction for all wave-vectors, with the maximum uncertainty estimated for $\epsilon(k)$.}
\label{tab:CvErr}
\end{table}
\endgroup


\newpage
\subsection{Normal fluid densities and relative densities at SVP}
\begingroup
\begin{table}[!h]
\resizebox{0.7\columnwidth}{!}{%
\begin{tabular}{|l|l|l|l|l|l|l|l|}
\hline
T    & $\rho$  & $\rho_n$   & $\rho_n^{Ph}$& $\rho_n^{R}$  & $\rho_n$/$\rho$  & $\rho_n^{Ph}$/$\rho$ & $\rho_n^{R}$/$\rho$ \\ \hline \hline
$K$    & $g/cm^3$& $g/cm^3$   &    $g/cm^3$    & $g/cm^3$        & --       & --        & --         \\ \hline
0.00	&	0.14514	&	--	&	--	&	--	&	--	&	--	&	--	\\
0.05	&	0.14514	&	1.1029E-10	&	1.1029E-10	&	3.9369E-74	&	7.59E-10	&	7.59E-10	&	2.71E-73	\\
0.10	&	0.14514	&	1.7539E-09	&	1.7539E-09	&	6.6529E-37	&	1.21E-08	&	1.20E-08	&	4.58E-36	\\
0.15	&	0.14514	&	8.7988E-09	&	8.7988E-09	&	1.5656E-24	&	6.06E-08	&	6.06E-08	&	1.07E-23	\\
0.20	&	0.14514	&	2.7495E-08	&	2.7495E-08	&	2.3054E-18	&	1.89E-07	&	1.89E-07	&	1.58E-17	\\
0.25	&	0.14514	&	6.6263E-08	&	6.6263E-08	&	1.1304E-14	&	4.56E-07	&	4.56E-07	&	7.78E-14	\\
0.30	&	0.14514	&	1.3550E-07	&	1.3550E-07	&	3.2126E-12	&	9.33E-07	&	9.33E-07	&	2.21E-11	\\
0.35	&	0.14514	&	2.4760E-07	&	2.4742E-07	&	1.7978E-10	&	1.70E-06	&	1.70E-06	&	1.23E-09	\\
0.40	&	0.14514	&	4.1961E-07	&	4.1596E-07	&	3.6495E-09	&	2.89E-06	&	2.86E-06	&	2.51E-08	\\
0.45	&	0.14514	&	6.9447E-07	&	6.5675E-07	&	3.7718E-08	&	4.78E-06	&	4.52E-06	&	2.59E-07	\\
0.50	&	0.14514	&	1.2304E-06	&	9.8722E-07	&	2.4320E-07	&	8.47E-06	&	6.80E-06	&	1.67E-06	\\
0.55	&	0.14514	&	2.5399E-06	&	1.4267E-06	&	1.1132E-06	&	1.75E-05	&	9.82E-06	&	7.67E-06	\\
0.60	&	0.14514	&	5.9393E-06	&	1.9968E-06	&	3.9426E-06	&	4.09E-05	&	1.37E-05	&	2.71E-05	\\
0.65	&	0.14513	&	1.4188E-05	&	2.7220E-06	&	1.1466E-05	&	9.77E-05	&	1.87E-05	&	7.90E-05	\\
0.70	&	0.14513	&	3.2202E-05	&	3.6307E-06	&	2.8571E-05	&	2.21E-04	&	2.50E-05	&	1.96E-04	\\
0.75	&	0.14513	&	6.7687E-05	&	4.7567E-06	&	6.2930E-05	&	4.66E-04	&	3.27E-05	&	4.33E-04	\\
0.80	&	0.14513	&	1.3156E-04	&	6.1411E-06	&	1.2542E-04	&	9.06E-04	&	4.23E-05	&	8.64E-04	\\
0.85	&	0.14513	&	2.3809E-04	&	7.8347E-06	&	2.3026E-04	&	0.00164	&	5.39E-05	&	0.00158	\\
0.90	&	0.14512	&	4.0476E-04	&	9.9013E-06	&	3.9486E-04	&	0.00278	&	6.82E-05	&	0.00272	\\
0.95	&	0.14512	&	6.5191E-04	&	1.2420E-05	&	6.3949E-04	&	0.00449	&	8.55E-05	&	0.00440	\\
1.00	&	0.14512	&	0.0010023	&	1.5487E-05	&	9.8677E-04	&	0.00690	&	1.06E-04	&	0.00679	\\
1.05	&	0.14512	&	0.0014804	&	1.9218E-05	&	0.0014612	&	0.0102	&	1.32E-04	&	0.0101	\\
1.10	&	0.14511	&	0.0021123	&	2.3749E-05	&	0.0020885	&	0.0145	&	1.63E-04	&	0.0143	\\
1.15	&	0.14511	&	0.0029246	&	2.9235E-05	&	0.0028954	&	0.0201	&	2.01E-04	&	0.0199	\\
1.20	&	0.14512	&	0.0039448	&	3.5848E-05	&	0.0039090	&	0.0271	&	2.47E-04	&	0.0269	\\
1.25	&	0.14512	&	0.0052006	&	4.3777E-05	&	0.0051569	&	0.0358	&	3.01E-04	&	0.0355	\\
1.30	&	0.14512	&	0.0067199	&	5.3222E-05	&	0.0066667	&	0.0463	&	3.66E-04	&	0.0459	\\ \hline
\end{tabular}%
}
\caption{Total density $\rho$ (from Ref. \onlinecite{DonnellyBarenghi}), normal fluid densities (total, phonon and roton contributions), and same quantities divided by the total density, as a function of temperature, calculated numerically using the measured dispersion relation, at saturated vapor pressure.}
\label{tab:rho_n}
\end{table}
\endgroup

\newpage
\section{Data files}
Data files provided for several parameters of interest, extending the tables present in the manuscript, are attached; their content is described  below. 

\begin{enumerate}
	\item DispersionP0allRange.txt  
	
	This file contains the data of the dispersion relation of $^4$He at saturated vapor pressure and very low temperatures (T$<$100\,mK). 
	For wave-vectors k$<$0.15\,\AA$^{-1}$: ultrasound data \cite{Rugar1984}. From 0.15 to 0.3\,\AA$^{-1}$: ultrasound and present neutron data combined within respective uncertainties; for k$>$ 0.3\,\AA$^{-1}$, present neutron data at  saturated vapor pressure and very low temperatures (T$<$0.1K). Details are given in the main article, together with an abstract of the present data table.
\begin{itemize}
\item First two lines are headers.
\item First column: wavevector in \AA$^{-1}$, from 0 to 3.6\,\AA$^{-1}$.
\item Second column: excitation energy in meV.
\item Third column: uncertainty of excitation energy in meV.
\end{itemize}

	\item DispersionAllPressures.txt
	
This file contains the data for several pressures, measured at low temperatures using an incident neutron energy of 3.52\,meV.
The initial raw data have been grouped in wave-vector bins of width 0.002\,\AA$^{-1}$, reducing the total number of data points from $\approx$90000 to $\approx$1000 points in the present file. 

Missing data points are indicated in the table as NAN, i.e. Not A Number; they correspond to a few gaps in the detector tubes set-up (see manuscript, section IV B).

The different columns correspond to (k, E, errE) at different pressures, where k is the wave-vector, E the excitation energy and errE the corresponding uncertainty.
The pressures (as indicated in the 3rd row of each column) are:
P= 0, 0.51, 1.02, 2.01, 5.01, 10.01, 24.08 bar
\begin{itemize}
\item First three lines are headers. The different columns are:
\item c1: row index
\item c2: Q (\AA$^{-1}$)
\item c3: E at P=0
\item c4: errE at P=0
\item c5: E at P=0.51 bar
\item c6: errE at P=0.51 bar
\item c7: E at P=1.02 bar
\item c8: errE at P=1.02 bar
\item c9: E at P=2.01 bar
\item c10: errE at P=2.01 bar
\item c11: E at P=5.01 bar
\item c12: errE at P=5.01 bar
\item c13: E at P=10.01 bar
\item c14: errE at P=10.01 bar
\item c15: E at P=24.08 bar
\item c16: errE at P=24.08 bar
\end{itemize}

	\item Cv-and-Entropy.txt
	
The file contains data of the thermodynamical properties at saturated vapor pressure, calculated using the measured dispersion curve. 
\begin{itemize}
\item First two lines are headers.  The different columns are:
\item c1: Temperature ($T$) in kelvin.
\item c2: Total specific heat ($C_V$)  in J/(K.mol). 
\item c3: uncertainty of $C_V$, calculated assuming a shift of the dispersion relation, in the same direction for all wave-vectors, with the maximum uncertainty estimated for $\epsilon(k)$. Above 0.8\,K the relative uncertainty of $C_V$ is almost constant at 1.1\%. Below this temperature, it decreases monotonically, reaching a minimum value of 0.13\%, deduced from the relative accuracy of the sound velocity data, which determine the very low temperature behavior.  
\item c4: phonon contribution $C_V^{Ph}$ (integral in the range $0<k<k_M$) in J/(K.mol). 
\item c5: roton contribution $C_V^{R}$ (integral in the range  $k_M<k<3.6$\,\AA$^{-1}$), 
\item c6: void
\item c7: total entropy $S$ in J/(K.mol).  
The maxon wave-vector at saturated vapor pressure is $k_M$=1.103\,\AA$^{-1}$.
\end{itemize}
	
  \item ExcitationsNumberDensity.txt

The number density of excitations (phonons, roton, and total number, per helium atom) as a function of temperature, at saturated vapor pressure. 
\begin{itemize}
\item First line contains headers.  The different columns are:
\item c1: Temperature ($T$) in kelvin.
\item c2: Number of phonons per helium atom.
\item c3: Number of rotons per helium atom.
\item c4: Number of excitations (phonons+rotons) per helium atom.
\end{itemize}

\end{enumerate}


\begin{thebibliography}{10}%
\makeatletter
\providecommand \@ifxundefined [1]{%
 \@ifx{#1\undefined}
}%
\providecommand \@ifnum [1]{%
 \ifnum #1\expandafter \@firstoftwo
 \else \expandafter \@secondoftwo
 \fi
}%
\providecommand \@ifx [1]{%
 \ifx #1\expandafter \@firstoftwo
 \else \expandafter \@secondoftwo
 \fi
}%
\providecommand \natexlab [1]{#1}%
\providecommand \enquote  [1]{``#1''}%
\providecommand \bibnamefont  [1]{#1}%
\providecommand \bibfnamefont [1]{#1}%
\providecommand \citenamefont [1]{#1}%
\providecommand \href@noop [0]{\@secondoftwo}%
\providecommand \href [0]{\begingroup \@sanitize@url \@href}%
\providecommand \@href[1]{\@@startlink{#1}\@@href}%
\providecommand \@@href[1]{\endgroup#1\@@endlink}%
\providecommand \@sanitize@url [0]{\catcode `\\12\catcode `\$12\catcode
  `\&12\catcode `\#12\catcode `\^12\catcode `\_12\catcode `\%12\relax}%
\providecommand \@@startlink[1]{}%
\providecommand \@@endlink[0]{}%
\providecommand \url  [0]{\begingroup\@sanitize@url \@url }%
\providecommand \@url [1]{\endgroup\@href {#1}{\urlprefix }}%
\providecommand \urlprefix  [0]{URL }%
\providecommand \Eprint [0]{\href }%
\providecommand \doibase [0]{http://dx.doi.org/}%
\providecommand \selectlanguage [0]{\@gobble}%
\providecommand \bibinfo  [0]{\@secondoftwo}%
\providecommand \bibfield  [0]{\@secondoftwo}%
\providecommand \translation [1]{[#1]}%
\providecommand \BibitemOpen [0]{}%
\providecommand \bibitemStop [0]{}%
\providecommand \bibitemNoStop [0]{.\EOS\space}%
\providecommand \EOS [0]{\spacefactor3000\relax}%
\providecommand \BibitemShut  [1]{\csname bibitem#1\endcsname}%
\let\auto@bib@innerbib\@empty
%</preamble>
\bibitem [{\citenamefont {Fetter}\ and\ \citenamefont
  {Walecka}(1971)}]{FetterWalecka}%
  \BibitemOpen
  \bibfield  {author} {\bibinfo {author} {\bibfnamefont {A.~L.}\ \bibnamefont
  {Fetter}}\ and\ \bibinfo {author} {\bibfnamefont {J.~D.}\ \bibnamefont
  {Walecka}},\ }\href@noop {} {\emph {\bibinfo {title} {{Q}uantum Theory of
  Many-Particle Systems}}}\ (\bibinfo  {publisher} {McGraw-Hill},\ \bibinfo
  {address} {New York},\ \bibinfo {year} {1971})\BibitemShut {NoStop}%
\bibitem [{\citenamefont {Jackson}\ \emph {et~al.}(1981)\citenamefont
  {Jackson}, \citenamefont {Jennings}, \citenamefont {Smith},\ and\
  \citenamefont {Lande}}]{SmithSolid}%
  \BibitemOpen
  \bibfield  {author} {\bibinfo {author} {\bibfnamefont {A.~D.}\ \bibnamefont
  {Jackson}}, \bibinfo {author} {\bibfnamefont {B.~K.}\ \bibnamefont
  {Jennings}}, \bibinfo {author} {\bibfnamefont {R.~A.}\ \bibnamefont {Smith}},
  \ and\ \bibinfo {author} {\bibfnamefont {A.}~\bibnamefont {Lande}},\
  }\href@noop {} {\bibfield  {journal} {\bibinfo  {journal} {Phys. Rev. B}\
  }\textbf {\bibinfo {volume} {24}},\ \bibinfo {pages} {105} (\bibinfo {year}
  {1981})}\BibitemShut {NoStop}%
\bibitem [{\citenamefont {Nozi\`eres}(2004)}]{NozSolid}%
  \BibitemOpen
  \bibfield  {author} {\bibinfo {author} {\bibfnamefont {P.}~\bibnamefont
  {Nozi\`eres}},\ }\href@noop {} {\bibfield  {journal} {\bibinfo  {journal} {J.
  Low Temp. Phys.}\ }\textbf {\bibinfo {volume} {137}},\ \bibinfo {pages} {45}
  (\bibinfo {year} {2004})}\BibitemShut {NoStop}%
\bibitem [{\citenamefont {Rugar}\ and\ \citenamefont
  {Foster}(1984)}]{Rugar1984}%
  \BibitemOpen
  \bibfield  {author} {\bibinfo {author} {\bibfnamefont {D.}~\bibnamefont
  {Rugar}}\ and\ \bibinfo {author} {\bibfnamefont {J.~S.}\ \bibnamefont
  {Foster}},\ }\href@noop {} {\bibfield  {journal} {\bibinfo  {journal} {Phys.
  Rev. B}\ }\textbf {\bibinfo {volume} {30}},\ \bibinfo {pages} {2595}
  (\bibinfo {year} {1984})}\BibitemShut {NoStop}%
\bibitem [{\citenamefont {Landau}(1941)}]{Landauroton}%
  \BibitemOpen
  \bibfield  {author} {\bibinfo {author} {\bibfnamefont {L.}~\bibnamefont
  {Landau}},\ }\href@noop {} {\bibfield  {journal} {\bibinfo  {journal} {J.
  Phys. (Moscow)}\ }\textbf {\bibinfo {volume} {5}},\ \bibinfo {pages} {71}
  (\bibinfo {year} {1941})}\BibitemShut {NoStop}%
\bibitem [{\citenamefont {Landau}(1947)}]{Landauroton2}%
  \BibitemOpen
  \bibfield  {author} {\bibinfo {author} {\bibfnamefont {L.}~\bibnamefont
  {Landau}},\ }\href@noop {} {\bibfield  {journal} {\bibinfo  {journal} {J.
  Phys. (Moscow)}\ }\textbf {\bibinfo {volume} {11}},\ \bibinfo {pages} {91}
  (\bibinfo {year} {1947})}\BibitemShut {NoStop}%
\bibitem [{\citenamefont {Stirling}(1983)}]{stirling-83}%
  \BibitemOpen
  \bibfield  {author} {\bibinfo {author} {\bibfnamefont {W.~G.}\ \bibnamefont
  {Stirling}},\ }in\ \href@noop {} {\emph {\bibinfo {booktitle} {{$75^{\rm
  th}$} Jubilee Conference on Helium-4}}},\ \bibinfo {editor} {edited by\
  \bibinfo {editor} {\bibfnamefont {J.~G.~M.}\ \bibnamefont {Armitage}}}\
  (\bibinfo  {publisher} {World Scientific},\ \bibinfo {address} {Singapore},\
  \bibinfo {year} {1983})\BibitemShut {NoStop}%
\bibitem [{\citenamefont {Glyde}(1994)}]{GlydeBook}%
  \BibitemOpen
  \bibfield  {author} {\bibinfo {author} {\bibfnamefont {H.~R.}\ \bibnamefont
  {Glyde}},\ }\href@noop {} {\emph {\bibinfo {title} {Excitations in Liquid and
  Solid Helium}}},\ Oxford Series on Neutron Scattering in Condensed Matter\
  (\bibinfo  {publisher} {Clarendon},\ \bibinfo {address} {Oxford},\ \bibinfo
  {year} {1994})\BibitemShut {NoStop}%
\bibitem [{\citenamefont {Stirling}(1991)}]{StirlingExeter}%
  \BibitemOpen
  \bibfield  {author} {\bibinfo {author} {\bibfnamefont {W.~G.}\ \bibnamefont
  {Stirling}},\ }in\ \href@noop {} {\emph {\bibinfo {booktitle} {Excitations in
  Two-Dimensional and Three-Dimensional {Q}uantum Fluids}}},\ \bibinfo {series}
  {NATO Advanced Study Institute, Series B: Physics}, Vol.\ \bibinfo {volume}
  {257},\ \bibinfo {editor} {edited by\ \bibinfo {editor} {\bibfnamefont
  {A.~F.~G.}\ \bibnamefont {Wyatt}}\ and\ \bibinfo {editor} {\bibfnamefont
  {H.~J.}\ \bibnamefont {Lauter}}}\ (\bibinfo  {publisher} {Plenum},\ \bibinfo
  {address} {New York},\ \bibinfo {year} {1991})\ pp.\ \bibinfo {pages}
  {25--46}\BibitemShut {NoStop}%
\bibitem [{\citenamefont {Donnelly}\ and\ \citenamefont
  {Barenghi}(1998)}]{DonnellyBarenghi}%
  \BibitemOpen
  \bibfield  {author} {\bibinfo {author} {\bibfnamefont {R.~J.}\ \bibnamefont
  {Donnelly}}\ and\ \bibinfo {author} {\bibfnamefont {C.~F.}\ \bibnamefont
  {Barenghi}},\ }\href {\doibase 10.1063/1.556028} {\bibfield  {journal}
  {\bibinfo  {journal} {J. Phys. Chem. Ref. Data}\ }\textbf {\bibinfo {volume}
  {27}},\ \bibinfo {pages} {1217} (\bibinfo {year} {1998})},\ \Eprint
  {http://arxiv.org/abs/https://doi.org/10.1063/1.556028}
  {https://doi.org/10.1063/1.556028} \BibitemShut {NoStop}%
\end{thebibliography}%


\end{document}

